\section{Methods}
\label{sec:methods}

The paraphrases are generated from an initial seed prompt that uses simplified language to
describe the desired results. The prompt represents a use-case in urban planning to find areas of
the city which have access to a nearby green space such as a park or a forest, called cool spot
analysis \citep{kluck2022}. It requires a multi-step process that includes querying feature types
within a geographic area, filtering on criteria, and then running a spatial processing tool on
the features. This use case is realistic and fairly simple to envision, consisting of just a few
steps. It has a clear correct answer that can only vary due to the selection of certain
parameters such as the area threshold (``Large enough for recreation'') or the number and size
of buffer intervals (``In incremental intervals''). While this does introduce some ambiguity, it
reflects a realistic prompt where not everything is specified precisely and the agent must
reason on its own to fill in the gaps. Since we have executed this workflow manually, we can use
its results as ground truth to validate the agent's output against it \citep{voigtlaender2022}.

\emph{Identify all areas in Münster that are within a 300-meter walking distance of a public park
or forest large enough for recreation. Output a visualization that highlights these well-served
areas and displays the walking distances from the green spaces in incremental intervals up to
the 300-meter limit.}

To create the dataset or paraphrases, a guided LLM approach was used. The paraphrase typology
was separated into individual functions which were used by an LLM to generate a paraphrase.
Then, an LLM is provided with distinct prompts corresponding to each paraphrase category. These
prompts incorporate specific instructions and contextual examples drawn from an established
paraphrase corpus \citep{kovatchev2018}.

To guarantee quality and relevance, every generated paraphrase undergoes a rigorous two-step
evaluation. First, an automated process runs basic checks: a stem check for morphological
paraphrases to ensure word stems match the original, and a cosine similarity check to filter out
near duplicates. Next, an ``LLM judge'' evaluates the output against predefined criteria, ensures
it is in the correct typological category (suggesting recategorization if needed), and checks
for grammatical correctness. In borderline cases where constraints are arguably met, the prompt
is flagged as ``relatively correct'' for manual human review. This dual-review process confirms
that semantic variations do not alter the core task requirements.

During the review phase, each paraphrase must satisfy eight distinct, hard constraints focused on
user intention. These strict guidelines are enforced to ensure that despite linguistic variations
or structural changes, the paraphrased prompt remains perfectly aligned with the original
geospatial request and expects the exact same practical outcome.

\begin{enumerate}
  \item \emph{The analysis is about the city of Münster.}
  \item \emph{The result is the areas that are within the distance limit of a qualifying green
    space.}
  \item \emph{The distance limit is 300 metres.}
  \item \emph{Distance is measured as walking distance (along walkable routes), not straight-line
    distance or by another mode.}
  \item \emph{The destinations are public parks or forests.}
  \item \emph{Only green spaces large enough for recreation count.}
  \item \emph{The output is a visualization that highlights the well-served areas.}
  \item \emph{The visualization shows walking distance from the green spaces in incremental
    intervals up to the limit.}
\end{enumerate}

\subsection{System Architecture}
\label{sec:architecture}

The core system is an agentic workflow built with LangGraph that translates natural language
requests into an executable workflow. This output workflow consists of a graph, not to be
confused with LangGraph, where the nodes are actionable technical steps and the edges represent
the data flow between steps. The agent processes the input prompt and from it builds a workflow
step by step, one node at a time. Each step will call a different sub-agent depending on the
stage of the current workflow graph and its most recent validation check.

The agent interfaces with an external geoserver to access the necessary tools for constructing
the workflow graphs. For this experimental setup, the computation server, pygeoapi
\citep{kralidis2026} exposes a controlled environment containing exactly 25~distinct geospatial
processes. Depending on the trial run, these processes are either exposed in full via the MCP
server or selectively retrieved using the RAG approach (Figure~\ref{fig:architecture}).

\begin{figure}[htbp]
  \caption{Overview of agentic architecture in test configuration with both RAG and MCP support}
  \label{fig:architecture}
  \includegraphics[width=0.7\textwidth, alt={Architecture diagram. Inside the LangGraph
    service, the agent workflow (LLM) receives a request from the web UI and returns a workflow
    graph. It exchanges information with a retriever, which reaches either a pgvector catalog via
    RAG or an MCP server via MCP. The web UI sends the graph to a runner for execution; the runner
    calls the OGC API Processes servers and returns the results to the web UI. The OGC API
    Processes servers are crawled to fill both the pgvector catalog and the MCP
    server.}]{01_architecture}
  \figsource{\ownwork}
\end{figure}

\subsection{Execution}
\label{sec:execution}

Testing the paraphrases required running each prompt individually through the LangChain process,
but the issue is that the primary interface is web-based. To circumvent this, queries are instead
performed directly within the LangGraph process with no network calls involved, isolating the
agentic interface to its core component. Every prompt is run separately with its own independent
context, and multiple prompts can be run concurrently. The entire workflow is created from start
to finish without interruption, and at each step of the agentic chain various metrics are
recorded. These include per-LLM call metrics such as the LangGraph node making the call, message
history, input/output tokens and cost, elapsed time, and visible tools. For each tool search, the
tools, servers, and closest matches are recorded. Total elapsed time and number of runs are saved
for each node. Finally, cumulative statistics are saved including: total time and iterations, the
nodes and unmet capabilities, LLM calls and token usage, and any errors that occurred. An
orchestration function manages all of the overhead processes needed to carry this work out.

While the batch execution process covered the baseline functionality of the experiment,
comparing the different tool retrieval methods required another layer of functionality to be
built on top of the existing test harness. This is accomplished through the use of a specialized
retrieval node in the LangGraph process. The retrieval node defined an abstract interface which
the different retrieval methods implemented in their concrete forms. The batch script can then be
configured to switch between these different variants of the retrieval node, resulting in an MCP
or RAG based agent. The main functionality of this specific node is to return a list of tools
from the server, either all tools in the case of MCP or the most similar tools in the case of RAG
(Figure~\ref{fig:batch}).

\begin{figure}[htbp]
  \caption{Overview of the experimental batch script configuration}
  \label{fig:batch}
  \includegraphics[width=\textwidth, alt={Batch setup diagram. Paraphrased prompts feed the
    agent workflow, which runs in-process without the web UI inside an orchestrator (batch
    script), with one context per concurrent run. A variant switch (MCP or RAG) configures the
    retriever, which queries either the MCP server for all tools or pgvector for the top-k
    similar tools. The runs write metrics (LLM calls, tool searches, per node and cumulative) to
    a store.}]{02_batch-setup}
  \figsource{\ownwork}
\end{figure}

\subsection{Evaluation}
\label{sec:evaluation}

The primary way that the workflow graphs were evaluated was through manual verification. Most of
the paraphrases created one of a small set of graphs that could be checked by hand. The graphs
were checked based on whether they would produce the desired results. However, one subset of the
graphs included extra nodes that were superfluous yet still produced an acceptable outcome. These
graphs were labeled ``leniently'' correct, while the exact match of the desired workflow was
labeled ``strictly'' correct (Figure~\ref{fig:graphs}).

The workflow graphs still needed to be checked against a known correct solution to ensure that
they worked correctly. An issue arises from the arbitrary values the agent chooses for some of
the geospatial functions. For example, a sufficiently large park is not defined by a specific
number in the prompt, leaving the agent to decide on its own. In some cases, the agent signaled
the intention to ask clarifying questions by labelling the value ``unknown.'' In these cases, a
reasonable default was used in the quantitative evaluation. Altogether, this means that a correct
solution can still vary from the ground truth by quite some degree. For this reason, quantitative
metrics of the output were only used for general validation rather than scoring.

\begin{figure}[htbp]
  \caption{Strictly correct (left) and leniently correct (right) workflow graphs}
  \label{fig:graphs}
  \includegraphics[height=7cm, alt={Strictly correct workflow graph with four nodes in a chain:
    city-boundary (place: Münster), osm-green-areas (parks, forests), min-area-filter (minArea)
    and multi-ring-buffer (rings up to 300 m).}]{03a_graph-strict}\qquad
  \includegraphics[height=7cm, alt={Leniently correct workflow graph: the same four nodes, plus
    a dashed, superfluous buffer node (300 m) that also takes the output of
    min-area-filter.}]{03b_graph-lenient}
  \figsource{\ownwork}
\end{figure}
